% main.tex - Documento raiz del paper de ingenieria / technical report.
% Variante IMRaD orientada a implementacion y evaluacion experimental.
% Compilar con pdflatex + bibtex.
\documentclass[11pt]{article}

\usepackage[utf8]{inputenc}
\usepackage[T1]{fontenc}
\usepackage[margin=2.5cm]{geometry}
\usepackage{booktabs}
\usepackage{hyperref}

% Marcadores academicos \todo{...} (ver .roo/rules/00_academic_core.md):
% fallback compilable. Si anades el paquete todonotes, elimina esta linea.
\providecommand{\todo}[1]{\textbf{[TODO: #1]}}

\title{TODO: titulo del paper}
\author{TODO: autor/es}
\date{TODO: fecha}

\begin{document}

\maketitle

\begin{abstract}
TODO: resumen de 150--250 palabras.
\end{abstract}

\input{src/01_introduction}
\input{src/02_related_work}
\input{src/03_architecture}
\input{src/04_implementation}
\input{src/05_evaluation}
\input{src/06_conclusion}

\appendix
\input{src/appendices/appendix_a}

\bibliographystyle{plain}
\bibliography{src/references}

\end{document}
